\documentclass[10pt,a4paper]{amsart}
\usepackage[margin=19mm]{geometry}
\usepackage{amsmath,amssymb,mathtools}
\usepackage[scr=rsfs,cal=euler]{mathalfa}
\usepackage{xfrac}
\usepackage[unicode,hidelinks,bookmarks=false]{hyperref}
\newcommand{\ie}{i.e.\ }
\newcommand{\cf}{cf.\ }
\newcommand{\resp}{respectively\ }
\newcommand{\Subsection}{\subsection}
\providecommand{\nobreakdash}{\nobreak\hbox{-}}
\setlength{\emergencystretch}{2em}
\multlinegap0pt
\usepackage{mathtools}
\usepackage{amssymb}
\usepackage{float}
\usepackage{xcolor}
\usepackage{tikz}
\newtheorem{definition}{Definition}
\newtheorem{lemma}{Lemma}
\newcommand{\bx}{\mathbf{x}}
\newcommand{\bz}{\mathbf{z}}
\title{On Stopping Rules and Spatial Adaptation for CART}
\author{Zineng Xu, Yuchao Cai,, Yan Shuo Tan}
\date{Source: arXiv:2608.15649v1 (CC BY 4.0)}
\begin{document}
\maketitle

\noindent\emph{Source and licence.} On Stopping Rules and Spatial Adaptation for CART. Version arXiv:2608.15649v1 (CC BY 4.0), distributed by arXiv under the Creative Commons Attribution 4.0 International licence, http://creativecommons.org/licenses/by/4.0/, as recorded in the arXiv licence metadata for this exact version and shown on the abstract page https://arxiv.org/abs/2608.15649v1. Full licence notice: \texttt{licenses/arxiv-2608.15649-CC-BY-4.0.txt}.\par\smallskip

\noindent\textit{Editorial scope.} Counted unit: the article's lemma lem:radius-one (source0.tex lines 5663--5692). The article's complete original proof of it is delivered with it (source0.tex lines 5694--5832, one \texttt{proof} environment of 139 lines). No mathematics is written, repaired or re-derived: the statement and the proof are the article's own lines, in the article's own notation, and no retained line is substituted or re-flowed.  Retained uncounted as the source-local context the proof needs: lines 915--926; lines 5646--5661.  Nothing was omitted from any retained span, so the package carries no omission record.  No retained line contains a citation, so the wrapper carries no bibliography and no bibliography entry.  No retained line contains a figure environment, an image-inclusion command or drawing code, so no image is delivered and the package declares no asset dependency.  Editorial formatting only, in addition: an AMS \texttt{amsart} preamble that restates the article's own declarations and macros used by the retained lines, namely the environments \texttt{definition}, \texttt{lemma} on separate counters, together with the standard packages those lines need.  The retained environments print in this document as Definition 1, Lemma 1 in order of appearance, whereas the article prints the same items with the numbering of its own full document.  The complete native source of the article as distributed in the arXiv e-print for this exact version is bundled unchanged as \texttt{sources/207770/source0.tex}.
\begin{definition}[Local anisotropic H\"older regularity]
\label{def::localholder}
Let $\alpha_j:[0,1]\to(0,1]$, $j\in[d]$, be measurable functions satisfying
$0<\alpha_{\min}\le\alpha_j(t)\le\alpha_{\max}\le1$ for all $j$ and $t$.
Let $L\ge1$ be a H\"older radius.
We say that a bounded regression function $f^*$ is locally anisotropic H\"older with exponents $\{\alpha_j\}_{j=1}^d$ and radius $L$ if, for all $\bx,\bx'\in[0,1]^d$,
\[
|f^*(\bx)-f^*(\bx')|
\le
L\sum_{j=1}^d |x_j-x_j'|^{\alpha_j(x_j)}.
\]
\end{definition}

\begin{definition}[Bounded local anisotropic H\"older regularity]
\label{def:local-holder-general}
Let $\alpha_j,\gamma_j:[0,1]\to(0,1]$, $j\in[d]$, be measurable
functions satisfying $0<\alpha_{\min}\le\alpha_j(t)\le\alpha_{\max}\le1$
for all $j$ and $t$. Let $L\ge1$ be a H\"older radius. We say that a
bounded regression function $f^*$ is \emph{locally anisotropic
H\"older with exponents $\{\alpha_j\}_{j=1}^d$, locality radii
$\{\gamma_j\}_{j=1}^d$, and radius $L$} if, for all
$\bx,\bx'\in[0,1]^d$ such that $|x_j-x_j'|<\gamma_j(x_j)$ for all
$j\in[d]$,
\begin{equation}\label{eq:general-local-holder}
|f^*(\bx)-f^*(\bx')|
\;\le\;
L\sum_{j=1}^d |x_j-x_j'|^{\alpha_j(x_j)} .
\end{equation}
\end{definition}

\begin{lemma}[The locality radii can be taken to be one]
\label{lem:radius-one}
Let $f^*$ satisfy Definition~\ref{def:local-holder-general} with
exponents $\{\alpha_j\}_{j=1}^d$, locality radii
$\{\gamma_j\}_{j=1}^d$, and radius $L$. Assume in addition that
\begin{enumerate}
\item[(i)] each $\alpha_j$ is H\"older continuous: there exist
  $\nu\in(0,1]$ and $L_\alpha\ge0$ such that
  $\lvert\alpha_j(t)-\alpha_j(t')\rvert\le
  L_\alpha\lvert t-t'\rvert^{\nu}$
  for all $t,t'\in[0,1]$ and $j\in[d]$;
\item[(ii)] the locality radii are uniformly bounded away from zero:
  there exists $c\in(0,1]$ such that $\gamma_j(t)\ge c$ for all
  $t\in[0,1]$ and $j\in[d]$.
\end{enumerate}
Then, for \emph{all} $\bx,\bx'\in[0,1]^d$,
\[
|f^*(\bx)-f^*(\bx')|
\;\le\;
L'\sum_{j=1}^d |x_j-x_j'|^{\alpha_j(x_j)},
\qquad
L' \;:=\; \Bigl(\frac{2}{c}\Bigr)^{1+L_\alpha} L .
\]
That is, $f^*$ is locally anisotropic H\"older with the same
exponents, radius $L'$, and locality radii $\gamma_j\equiv 1$; in
particular, it satisfies the unrestricted form of the local H\"older
condition used in the main analysis, with $L$ replaced by $L'$.
% Point the final clause at the main-text definition, e.g.
% Definition~\ref{def::localholder}, once labels are wired up.
\end{lemma}

\begin{proof}
The proof consists of a reduction to one dimension followed by a
chaining argument; assumption (i) enters only through the two
exponent-comparison displays in the proof of the Claim below, and
assumption (ii) only through the admissibility of the chain steps and
the bound $m\le 2/c$ on their number.

For the one-dimensional reduction we write, for measurable functions
$\alpha:[0,1]\to(0,1]$ and $\gamma:[0,1]\to(0,1]$ and a constant
$M>0$,
\begin{equation}\label{eq:1d-class}
\begin{aligned}
\mathcal{H}_1(\alpha,\gamma,M)
\coloneqq\bigl\{g:[0,1]\to\mathbb{R}:\;&
\lvert g(t')-g(t)\rvert
\le M\lvert t-t'\rvert^{\alpha(t)}\\
&\text{for all $t,t'\in[0,1]$ with
$\lvert t-t'\rvert<\gamma(t)$}\bigr\}.
\end{aligned}
\end{equation}

\medskip
\noindent\textit{Step 1: Reduction to one dimension.} Fix
$\bx,\bx'\in[0,1]^d$ and define the coordinate path
$\bz^{(0)}:=\bx$ and
\[
  \bz^{(j)} \;:=\; (x_1',\ldots,x_j',\,x_{j+1},\ldots,x_d),
  \qquad j\in[d],
\]
so that $\bz^{(d)}=\bx'$ and consecutive points differ only in one
coordinate. By the triangle inequality,
\begin{equation}\label{eq:coordinate-path}
  |f^*(\bx)-f^*(\bx')|
  \;\le\; \sum_{j=1}^d
  \bigl|f^*(\bz^{(j-1)})-f^*(\bz^{(j)})\bigr| .
\end{equation}
For each $j$, define the univariate section
$g_j(t):=f^*(x_1',\ldots,x_{j-1}',\,t,\,x_{j+1},\ldots,x_d)$,
$t\in[0,1]$, so that the $j$th summand in \eqref{eq:coordinate-path}
equals $|g_j(x_j)-g_j(x_j')|$. We first record that
$g_j\in\mathcal{H}_{1}(\alpha_j,\gamma_j,L)$: if
$t,t'\in[0,1]$ satisfy $|t-t'|<\gamma_j(t)$, the two points of
$[0,1]^d$ at which $g_j(t)$ and $g_j(t')$ evaluate $f^*$ agree in
every coordinate $k\neq j$, so the admissibility condition
$|z_k-z_k'|=0<\gamma_k(z_k)$ holds off coordinate $j$, and
\eqref{eq:general-local-holder} yields
$|g_j(t')-g_j(t)|\le L\,|t-t'|^{\alpha_j(t)}$, the terms with
$k\neq j$ vanishing. It therefore suffices to prove the following
one-dimensional claim; applying it to $g_j$ with base point $x_j$ and
increment $x_j'-x_j$, and summing \eqref{eq:coordinate-path} over $j$,
gives the lemma with the dimension-free constant
$L'=(2/c)^{1+L_\alpha}L$.

\medskip
\noindent\emph{Claim.} Let
$g\in\mathcal{H}_{1}(\alpha,\gamma,L)$ with
$\gamma(\cdot)\ge c$ on $[0,1]$ and
$\lvert\alpha(t)-\alpha(t')\rvert\le L_\alpha\lvert t-t'\rvert^{\nu}$
for all $t,t'\in[0,1]$. Then, for every $x\in[0,1]$ and every $h\neq0$
with $x+h\in[0,1]$,
\[
  \lvert g(x+h)-g(x)\rvert
  \;\le\; \Bigl(\frac{2}{c}\Bigr)^{1+L_\alpha} L\,
  \lvert h\rvert^{\alpha(x)} .
\]

\medskip
\noindent\emph{Proof of the Claim.} If $\lvert h\rvert<\gamma(x)$, the
bound holds with constant $L\le (2/c)^{1+L_\alpha}L$, directly from
\eqref{eq:1d-class}. Assume therefore $\lvert h\rvert\ge\gamma(x)$.
Without loss of generality, we assume $h>0$: the reflected function
$\tilde g(u):=g(1-u)$ belongs to
$\mathcal{H}_1\bigl(\alpha(1-\cdot),\gamma(1-\cdot),L\bigr)$,
whose exponent and radius functions satisfy (i) and (ii) with the same
$(\nu,L_\alpha,c)$, and
$g(x+h)-g(x)=\tilde g\bigl((1-x)+\lvert h\rvert\bigr)-\tilde g(1-x)$
with $\alpha\bigl(1-(1-x)\bigr)=\alpha(x)$, so the case $h<0$ follows
from the case $h>0$ applied to $\tilde g$.

Let $m=\lfloor 2h/c\rfloor$. Since $c\le\gamma(x)\le h\le 1$, we have
$2\le m\le 2/c$. Let $x_i:=x+ic/2$ for $0\le i\le m-1$; from
$m\le 2h/c<m+1$, we have
\[
  c/2 \;\le\; x+h-x_{m-1} \;<\; c .
\]
By the triangle inequality, we have
\begin{align*}
|g(x+h)-g(x)|&\leq \sum_{i=1}^{m-1} |g(x_i)-g(x_{i-1})|
  + |g(x+h)-g(x_{m-1})|.
\end{align*}
The local
H\"older condition \eqref{eq:1d-class} gives
\begin{align*}
|g(x+h)-g(x)|\leq L\sum_{i=1}^{m-1}|x_i-x_{i-1}|^{\alpha(x_i)}
  + L|x+h-x_{m-1}|^{\alpha(x+h)}.
\end{align*}
Using the H\"older condition of $\alpha$ in (i), we have
\begin{align*}
|x_i-x_{i-1}|^{\alpha(x_i)}\leq
|x_i-x_{i-1}|^{\alpha(x)-L_\alpha|x_i-x|^{\nu}}
\end{align*}
and
\begin{align*}
|x+h-x_{m-1}|^{\alpha(x+h)}\leq
|x+h-x_{m-1}|^{\alpha(x)-L_\alpha h^{\nu}}.
\end{align*}
Since $|x_i-x|^{\nu}\le 1$ and $h^{\nu}\le 1$, while every increment
lies in $[c/2,1]$, we have
\begin{align*}
|x_i-x_{i-1}|^{L_\alpha|x_i-x|^{\nu}}
\;\geq\; |x_i-x_{i-1}|^{L_\alpha}
\;\geq\; (c/2)^{L_\alpha}
\end{align*}
and
\begin{align*}
|x+h-x_{m-1}|^{L_\alpha h^{\nu}}
\;\geq\; |x+h-x_{m-1}|^{L_\alpha}
\;\geq\; (c/2)^{L_\alpha}.
\end{align*}
Therefore, since every increment is at most $c\le h$ and there are
$m\le 2/c$ of them, we get
\begin{align*}
|g(x+h)-g(x)|& \leq
\Bigl(\frac{2}{c}\Bigr)^{L_\alpha}
L\sum_{i=1}^{m-1}|x_i-x_{i-1}|^{\alpha(x)}
  + \Bigl(\frac{2}{c}\Bigr)^{L_\alpha} L\,
  |x+h-x_{m-1}|^{\alpha(x)}\\
& \leq \Bigl(\frac{2}{c}\Bigr)^{L_\alpha} L m\, h^{\alpha(x)}
\;\leq\; \Bigl(\frac{2}{c}\Bigr)^{1+L_\alpha} L\, h^{\alpha(x)}.
\end{align*}
This proves the Claim.

\medskip
Finally, since $|x_j-x_j'|\le 1$ and $\alpha_j(x_j)>0$ imply
$|x_j-x_j'|^{\alpha_j(x_j)}\le1$, the conclusion of the lemma also
gives $\sup_{\bx,\bx'}|f^*(\bx)-f^*(\bx')|\le L'd$, so the boundedness
postulated in Definition~\ref{def:local-holder-general} holds
automatically.
\end{proof}

\end{document}
